\section{Related work}
\label{sec:related}

\paragraph{Elevator traffic analysis.} The engineering of lifts has a mature
literature on round-trip time, handling capacity and interval --- how many
cars a building needs and how long people wait --- codified in CIBSE Guide D
\cite{cibse2020guide} and the treatments of Barney and Al-Sharif
\cite{barney2015elevator} and Strakosch \cite{strakosch2010vertical}. Those
analyses take the population's floor distribution as given and size the
system; we hold the system fixed and ask the passenger's question. The
kinematic model (rated speed, acceleration, door and transfer times) and the
collective-control logic we simulate come from that literature.

\paragraph{Facility location on a line.} Minimising expected distance to a
random point on a line is the 1-median problem, solved by the median of the
destination distribution \cite{hakimi1964optimum}; a fixed ``retrieval''
cost that depends on where a server is parked turns it into a location problem
with a mobile server, close in spirit to the analysis of where to park a
single ambulance or repair vehicle \cite{larson1974hypercube}. Our closed
form is the 1-median with the elevator's parking position added to the
distance.

\paragraph{Stairs versus elevator.} Studies of stair use in buildings measure
walking speeds per flight and the reluctance to climb more than a few floors
\cite{fujiyama2010ascending}; the asymmetry between ascending and descending
speeds is well documented and is what motivated our stairs model.

\paragraph{Simulation of elevator systems.} Discrete-event simulation is the
standard tool for evaluating dispatch and group control
\cite{siikonen1993elevator}; our simulator is a minimal member of that family
(collective control, nearest-car dispatch, capacity), chosen so that it can be
validated exactly against the closed form at zero traffic and so that every
effect we report can be traced to a single mechanism.
